
This paper investigated whether a language model's architecture family can be identified from its adversarial failures alone. A large architecture-family effect is observed --- $\eta^{2}=0.293$ (Cohen's $f^{2}\approx 0.41)$ --- consistent across 83\% of attack categories at base model scale. The strongest differentiation occurs on adversarial NLI entailment tasks (adv\_rte, $\eta^{2}=0.592)$, consistent with the hypothesis that bidirectional and causal attention mechanisms diverge on tasks requiring global semantic integration. Statistical significance at $\alpha =0.05$ is not achieved (p=0.147), which is expected given that N=9 models provide approximately 40\% power at this effect size. Leave-one-model-out classification at N=3 per family is geometrically degenerate and achieves chance accuracy, a result that does not bear on the existence of family structure. A validated seven-module pipeline (1,788 lines) supporting replication and extension is provided.

A confirmatory study (h-e1-v2) with $N\geq 15$ models ($\geq 5$ per family) is the direct next step: power analysis projects approximately 80\% power at the observed effect size. The encoder-decoder family requires MNLI fine-tuning for T5 and BART to resolve the ANLI-R3 coverage gap. After confirmation, the planned mechanism chain --- attention-concentration analysis ($\Delta C$ extraction) and mediation testing --- can proceed from the validated existence result reported here.

Future directions include: the h-e1-v2 confirmatory expansion; alternative classifiers (LDA, centroid) tested on the existing N=9 data; CheckList behavioral partition; surrogate-bias decomposition across AdvGLUE, ANLI-R3, and CheckList; targeted entailment study using SNLI given the adv\_rte peak finding; and extension to large-scale (7B+) models.

